\documentclass[10pt,a4paper]{amsart}
\usepackage[margin=19mm]{geometry}
\usepackage{amsmath,amssymb,mathtools}
\usepackage[scr=rsfs,cal=euler]{mathalfa}
\usepackage{xfrac}
\usepackage[unicode,hidelinks,bookmarks=false]{hyperref}
\newcommand{\ie}{i.e.\ }
\newcommand{\cf}{cf.\ }
\newcommand{\resp}{respectively\ }
\newcommand{\Subsection}{\subsection}
\providecommand{\nobreakdash}{\nobreak\hbox{-}}
\setlength{\emergencystretch}{2em}
\multlinegap0pt
\usepackage{amssymb}
\newtheorem{proposition}{Proposition}
\title{Direct Generation of a Somos--4 Sequence\\ from an Algebraic Generating Function}
\author{Thomas Scheuerle}
\date{Source: arXiv:2609.15754v1 (CC BY 4.0)}
\begin{document}
\maketitle

\noindent\emph{Source and licence.} Direct Generation of a Somos--4 Sequence\\ from an Algebraic Generating Function. Version arXiv:2609.15754v1 (CC BY 4.0), distributed by arXiv under the Creative Commons Attribution 4.0 International licence, http://creativecommons.org/licenses/by/4.0/, as recorded in the arXiv licence metadata for this exact version and shown on the abstract page https://arxiv.org/abs/2609.15754v1. Full licence notice: \texttt{licenses/arxiv-2609.15754-CC-BY-4.0.txt}.\par\smallskip

\noindent\textit{Editorial scope.} Counted unit: the article's proposition prop:tail-quadratic (source0.tex lines 307--324). The article's complete original proof of it is delivered with it (source0.tex lines 326--362, one \texttt{proof} environment of 37 lines). No mathematics is written, repaired or re-derived: the statement and the proof are the article's own lines, in the article's own notation, and no retained line is substituted or re-flowed.  Retained uncounted as the source-local context the proof needs: lines 287--291; lines 299--303.  Nothing was omitted from any retained span, so the package carries no omission record.  No retained line contains a citation, so the wrapper carries no bibliography and no bibliography entry.  No retained line contains a figure environment, an image-inclusion command or drawing code, so no image is delivered and the package declares no asset dependency.  Editorial formatting only, in addition: an AMS \texttt{amsart} preamble that restates the article's own declarations and macros used by the retained lines, namely the environments \texttt{proposition} on separate counters, together with the standard packages those lines need.  The retained environments print in this document as Proposition 1 in order of appearance, whereas the article prints the same items with the numbering of its own full document.  The complete native source of the article as distributed in the arXiv e-print for this exact version is bundled unchanged as \texttt{sources/210446/source0.tex}.
\begin{equation}
 T_n(x)=\cfrac{1}{1-\cfrac{d_nx}{1-\cfrac{d_{n+1}x}{1-\ddots}}},
 \qquad T_n=\frac{1}{1-d_nxT_{n+1}}.
 \label{eq:tail-fraction}
\end{equation}

\begin{equation}
 d_n=\frac1{d_{n-1}d_{n-2}}\quad(n\geq3\text{ odd}),\qquad
 d_n=\frac{t}{d_{n-2}}+d_{n-3}\quad(n\geq4\text{ even}).
 \label{eq:isolated-tail-recurrence}
\end{equation}

\begin{proposition}[Quadratic equation for the recurrent tail]
\label{prop:tail-quadratic}
Let $T=T_1$ be the formal continued fraction defined by
\eqref{eq:tail-fraction}--\eqref{eq:isolated-tail-recurrence}. Then $T(0)=1$ and
\begin{equation}
 \mathcal A(x)T^2+\mathcal B(x)T+\mathcal C(x)=0,
 \label{eq:tail-quadratic}
\end{equation}
where
\begin{align}
 \mathcal A(x)&=(bt+1)x-b\bigl((b+d)t+1\bigr)x^2,
 \label{eq:tail-A}\\
 \mathcal B(x)&=bdt\,x^2+\bigl(-bt+b^2d+d^2b-1\bigr)x-bd,
 \label{eq:tail-B}\\
 \mathcal C(x)&=bd-bd^2x.
 \label{eq:tail-C}
\end{align}
\end{proposition}

\begin{proof}
Removing one level of a Stieltjes fraction is a fractional-linear
transformation,
\begin{equation}
 T_{n+1}=\frac{T_n-1}{d_nxT_n}.
 \label{eq:tail-extraction}
\end{equation}
Consequently, quadratic relations are preserved under tail extraction. More
explicitly, if $Q_n(Y)$ is quadratic and $Q_n(T_n)=0$, then a polynomial for
the next tail is
\begin{equation}
 Q_{n+1}(Y)=(1-d_nxY)^2
 Q_n\!\left(\frac1{1-d_nxY}\right),
 \label{eq:quadratic-transport}
\end{equation}
For the quadratic in \eqref{eq:tail-quadratic}, write
$Q_{b,d}(x,Y)=\mathcal A(x)Y^2+\mathcal B(x)Y+\mathcal C(x)$ and set
\[
 b'=\frac1{bd},\qquad d'=b+\frac{t}{d}.
\]
A direct expansion gives the exact covariance identity
\begin{equation}
 (1-bx-dxZ)^2Q_{b,d}\!\left(x,\frac{1-dxZ}{1-bx-dxZ}\right)
 =\frac{b^3d^4x^2}{bd+t}\,Q_{1/(bd),\,b+t/d}(x,Z).
 \label{eq:two-tail-covariance}
\end{equation}
For the recurrent tails,
\[
 T_1=\frac{1-dxT_3}{1-bx-dxT_3}.
\]
The recurrence \eqref{eq:isolated-tail-recurrence} maps the parameter pair
$(b,d)$ for $T_1$ to $(1/(bd),b+t/d)$ for $T_3$.  Iterating
\eqref{eq:two-tail-covariance} therefore shows, for every $N$, that
$Q_{b,d}(x,T_1)$ is divisible by $x^{2N}$ in
$\mathbb Q(b,d,t)[[x]]$.  Hence it vanishes identically.  Finally,
$Q_{b,d}(0,Y)=bd(1-Y)$, so the formal branch satisfying $T_1(0)=1$ is unique.
\end{proof}

\end{document}
